\section{El nuevo panorama de la revisión de código con IA}

\subsection{Herramientas principales}

\begin{itemize}
    \item \textbf{Graphite}: el invitado de esta semana, Tomas Reimers, es su CPO
    \item \textbf{Greptile}: revisión con IA basada en la comprensión del repositorio
    \item \textbf{CodeRabbit}: revisión automatizada de PR
    \item \textbf{Claude Code / Codex}: herramientas de codificación con IA de propósito general que también sirven para revisión
\end{itemize}

\subsection{Los cambios que trae la revisión de código con IA}

\begin{enumerate}
    \item \textbf{Mayor eficiencia}: automatiza las verificaciones rutinarias y acorta el ciclo de revisión
    \item \textbf{Consistencia}: la IA aplica los mismos estándares a todo el código
    \item \textbf{Compartición de conocimiento}: la IA puede portar el contexto de todo el repositorio
    \item \textbf{Menor carga cognitiva}: permite que las personas se concentren en la lógica compleja y las decisiones de arquitectura
    \item \textbf{Mejora continua}: los sistemas de IA acumulan más datos con el tiempo y mejoran
    \item \textbf{Comprensión integral}: la IA puede examinar el impacto de un cambio desde una perspectiva global
\end{enumerate}

\begin{knowledgebox}{El reparto de tareas en la revisión de código con IA}
La mejor práctica no es que la IA sustituya por completo a la revisión humana, sino que la IA se encargue de los problemas ``de superficie'' (formato, nomenclatura, patrones de bugs comunes, verificaciones de seguridad) y las personas se concentren en los problemas ``de fondo'' (idoneidad de la arquitectura, corrección de la lógica de negocio, mantenibilidad a largo plazo).
\end{knowledgebox}

\subsection{Las limitaciones actuales}

\begin{itemize}
    \item \textbf{Costo de configuración}: requiere una inversión inicial para entrenar el sistema
    \item \textbf{Falsos positivos}: es necesario entrenar el sistema continuamente para reducirlos
    \item \textbf{No capta los usos y las mejores prácticas propios del repositorio}: los idiom propios de cada repositorio siguen estando fuera del alcance de la IA
    \item \textbf{No puede manejar lógica de negocio ni decisiones de arquitectura complejas}: es justo donde las personas siguen siendo indispensables
    \item \textbf{Los cambios de seguridad exigen especial cautela}: la IA puede pasar por alto implicaciones de seguridad sutiles
    \item \textbf{Suele pasar por alto casos extremos}
\end{itemize}

\begin{warningbox}{En la era de la IA, la revisión de código es más importante, no menos}
Cuando la IA escribe la mayor parte de tu código, la revisión de código se vuelve \textbf{más} importante. Tú eres responsable del código que se fusiona y se despliega: no puedes echarle la culpa a la IA. ``The AI wrote it'' no es una excusa aceptable.
\end{warningbox}

\subsection{La arquitectura de implementación de un sistema de revisión con IA}

Si se piensa en la revisión con IA como un sistema de ingeniería y no como un chatbot, incluye al menos el siguiente flujo:

\begin{table}[H]
\centering
\begin{tabular}{p{0.18\textwidth} p{0.25\textwidth} p{0.4\textwidth}}
\toprule
Etapa & Entradas clave & Tarea principal \\
\midrule
Análisis del cambio & Git diff, tipo de archivo, mensaje de commit & Determinar qué módulos merecen atención y qué cambios tienen alto riesgo \\
Recuperación del repositorio & Archivos relevantes, PR históricos, normas de código & Complementar el contexto que el diff no cubre, para reducir comentarios fuera de contexto \\
Ejecución de reglas & Lint, reglas de seguridad, restricciones del equipo & Filtrar los problemas evidentes y reducir la carga humana \\
Evaluación semántica & Flujo lógico, contratos de interfaz, rutas de excepción & Detectar riesgos más profundos de corrección o mantenibilidad \\
Generación de comentarios & Nivel de riesgo, evidencia, propuesta de solución & Producir review comments ejecutables, discutibles y atribuibles \\
\bottomrule
\end{tabular}
\caption{El flujo típico de la revisión de código con IA}
\end{table}

\begin{warningbox}{El modo de falla más grande es el ensamblaje incorrecto del contexto}
El problema de muchos productos de revisión con IA no es que el modelo no sepa redactar comentarios, sino que no se le suministra suficiente contexto del repositorio, reglas del equipo ni convenciones históricas. En cuanto el contexto se desalinea, se generan muchos comentarios que parecen razonables pero que en realidad no sirven.
\end{warningbox}

\subsection{Dónde el reviewer humano es más insustituible}

Tanto el curso como la práctica de la industria apuntan a la misma conclusión: cuanto mejor sea la IA para revisar patrones locales, más debe el reviewer humano dedicar su tiempo al juicio de alto nivel.

\begin{itemize}
    \item \textbf{Verificación de la intención}: ¿este cambio de verdad resuelve el problema correcto?
    \item \textbf{Coherencia de la arquitectura}: ¿hace que el sistema evolucione hacia más desorden?
    \item \textbf{Memoria organizacional}: ¿viola alguna regla implícita que el equipo formó tras experiencias pasadas?
    \item \textbf{Compensación de riesgos}: ¿algunos defectos son aceptables?, ¿alguna deuda técnica debería posponerse?
\end{itemize}

\begin{importantbox}{El reviewer del futuro se parece más a un ``diseñador de sistemas que firma al final''}
A medida que aumenta la proporción de código generado por IA, el reviewer ya no es solo quien revisa el formato o la nomenclatura, sino quien responde por si un cambio debe entrar a la rama principal. Este papel exige una comprensión de sistemas más sólida, no una lectura más rápida del diff.
\end{importantbox}

\subsection{Cómo debería un equipo medir si su sistema de revisión de verdad mejoró}

Sin métricas, los equipos suelen confundir ``más cantidad de comentarios'' con ``mayor calidad de revisión''. Un indicador más razonable debe cubrir a la vez velocidad, calidad y confianza.

\begin{itemize}
    \item \textbf{Tiempo de la primera respuesta}: cuánto tarda el autor en recibir el primer lote de comentarios útiles tras enviar el cambio
    \item \textbf{Tasa de defectos que escapan}: proporción de problemas que llegan a producción después de fusionarse
    \item \textbf{Tasa de falsos positivos}: proporción de comentarios de la IA que resultan inválidos
    \item \textbf{Proporción de PR reabiertos}: indica si la revisión inicial de verdad detectó los problemas centrales
    \item \textbf{Carga del reviewer}: cuánto tiempo consumen los ingenieros experimentados en comentarios de bajo valor
\end{itemize}

\begin{knowledgebox}{Un buen sistema de revisión no hace más ruidoso cada PR, sino más claro el PR de alto riesgo}
Si la IA genera muchos comentarios plantilla en cada diff, el equipo pronto desarrollará el hábito de ignorarlos. Un sistema de verdad eficaz debe concentrar la atención en los problemas más críticos, para que los comentarios importantes se atiendan con más facilidad.
\end{knowledgebox}

\subsection{Resumen de la sección}

Las herramientas de revisión de código con IA elevan de forma notable la eficiencia y la consistencia al automatizar las verificaciones rutinarias, pero la lógica de negocio compleja, las decisiones de arquitectura y las mejores prácticas propias de cada repositorio todavía requieren revisores humanos. En la era de la codificación con IA, la importancia de la revisión de código no disminuye, sino que aumenta.
